\documentclass[11pt,a4paper]{article}

% -----------------------------------------------------------------------------
% Paquetes: todos forman parte de una instalacion estandar de LaTeX.
% -----------------------------------------------------------------------------
\usepackage[utf8]{inputenc}
\usepackage[T1]{fontenc}
\usepackage{enumitem}
% Fecha en espanol sin depender de babel ni de spanish.ldf.
\newcommand{\FechaEnEspanol}{%
  \number\day\space de\space
  \ifcase\month
    \or enero\or febrero\or marzo\or abril\or mayo\or junio\or
    julio\or agosto\or septiembre\or octubre\or noviembre\or diciembre
  \fi\space de\space\number\year
}
\usepackage[a4paper,margin=2.5cm,headheight=42pt]{geometry}
\usepackage{graphicx}
\usepackage{xcolor}
\usepackage{fancyhdr}

% lastpage no es necesario para compilar, pero permite mostrar "pagina de N".
\newif\ifHayUltimaPagina
\IfFileExists{lastpage.sty}{
  \usepackage{lastpage}
  \HayUltimaPaginatrue
}{
  \HayUltimaPaginafalse
}

% -----------------------------------------------------------------------------
% CONFIGURACION DEL TRABAJO: editar solamente este bloque para reutilizarlo.
% -----------------------------------------------------------------------------
\newcommand{\NombreMateria}{Redes de Computadoras}
\newcommand{\TipoTrabajo}{Trabajo Práctico}
\newcommand{\NumeroTrabajo}{2}
% Cambiar por "Teórico" o "Práctico" según corresponda.
\newcommand{\TipoTP}{Práctico}
\newcommand{\NombreTrabajo}{Conceptos de capa física y capa de enlace de datos.}
\newcommand{\NombreGrupo}{NetRunners}
\newcommand{\Docente}{SANTIAGO MARTIN HENN}
\newcommand{\LogoArchivo}{logo.png}
\newcommand{\LogoRuta}{\LogoArchivo}

% Una linea por integrante. Se pueden agregar o quitar filas.
\newcommand{\Integrantes}{%
  \begin{tabular}{ll}
    Bastida, Lucas Ramiro & DNI 39444511 \\
    Contessi, Ruben Andres & DNI 33315235 \\
    Garay, Ignacio Hernan & DNI 44191925 \\
    Giraudo, Juan Pablo & DNI 34970089 \\
    Peñaloza, Gonzalo Adrian & DNI 40441355 \\
    Quinteros del Castillo, Tomás Agustín & DNI 46169739 \\
    Vasconsellos Blason, Benjamin Alejandro & DNI 45642879 \\
  \end{tabular}%
}

% -----------------------------------------------------------------------------
% Elementos comunes de caratula y encabezado.
% -----------------------------------------------------------------------------
\newcommand{\InsertarLogo}[1][0.24\textwidth]{%
  \ifx\LogoRuta\empty
    \fbox{\parbox[c][1cm][c]{#1}{\centering Logo institucional\\\small (colocar archivo)}}%
  \else
    \includegraphics[width=#1,keepaspectratio]{\LogoRuta}%
  \fi
}

\newcommand{\TituloCorto}{TP \NumeroTrabajo (\TipoTP): \NombreTrabajo}

\pagestyle{fancy}
\fancyhf{}
\fancyhead[L]{\InsertarLogo[1.4cm]}
\fancyhead[C]{\small\textbf{\TituloCorto}}
\fancyhead[R]{\small\NombreGrupo}
\renewcommand{\headrulewidth}{0.4pt}
\fancyfoot[C]{\small Pagina \thepage\ifHayUltimaPagina\ de \pageref{LastPage}\fi}
\renewcommand{\footrulewidth}{0pt}

\begin{document}

\begin{titlepage}
  \thispagestyle{empty}
  \begin{center}
    \InsertarLogo[0.32\textwidth]

    \vfill
    {\Large\textsc{\NombreMateria}\par}
    \vspace{1.5cm}
    {\Huge\bfseries \TipoTrabajo\par}
    \vspace{0.35cm}
    {\LARGE N.\textsuperscript{o} \NumeroTrabajo (\TipoTP)\par}
    \vspace{0.8cm}
    {\Large\NombreTrabajo\par}

    \vfill
    \begin{tabular}{rl}
      \textbf{Grupo:} & \NombreGrupo \\
      \textbf{Docente:} & \Docente
    \end{tabular}

    \vspace{1cm}
    \begin{center}
      \textbf{Integrantes}\\[0.2cm]
      \Integrantes
    \end{center}

    \vfill
    {\large \FechaEnEspanol\par}
  \end{center}
\end{titlepage}

\section{Consignas}
\begin{enumerate}[label=\arabic*)]
  \item
  \begin{enumerate}[label=\alph*)]
    \item
    El fenómeno físico representado en la figura es el \textbf{ruido},
    particularmente el \textbf{ruido impulsivo}. El ruido consiste en señales
    no deseadas que se insertan en algún punto entre el emisor y el receptor,
    modificando la señal transmitida.

    En la figura se observa una señal sinusoidal que presenta una perturbación
    en una parte de la transmisión. Esta perturbación está formada por
    variaciones irregulares y de corta duración, características del ruido
    impulsivo.

    El ruido impulsivo es de carácter no continuo y está compuesto por pulsos
    o picos irregulares de corta duración y amplitud relativamente grande.
    Este tipo de ruido puede provocar errores en los datos recibidos,
    siendo especialmente importante en las comunicaciones digitales.
    \item
    \item
  \end{enumerate}

  \item
  \item
  \begin{enumerate}[label=\alph*)]
    \item
    \item
    \item
  \end{enumerate}

  \item

  \item
  \begin{enumerate}[label=\alph*)]
    \item
    \item
    \item
    \item
  \end{enumerate}

  \item
  \begin{enumerate}[label=\alph*)]
    \item
    \item
  \end{enumerate}
\end{enumerate}

\end{document}
